\section{Metodyka badawcza}
\label{sec:metodyka}

\subsection{Ogólny schemat eksperymentu}

Aby porównanie było uczciwe, wszystkie modele uczono i testowano \emph{na tych
samych danych} i oceniano \emph{tymi samymi miarami}. Cały eksperyment przebiega
według prostego scenariusza:

\begin{enumerate}
  \item wczytanie danych (rzeczywisty zbiór CSE-CIC-IDS2018 lub --- gdy nie jest
        dostępny --- dane syntetyczne),
  \item przygotowanie danych (oczyszczenie, ujednolicenie, podział na zbiory),
  \item wytrenowanie pięciu modeli: regresji logistycznej, lasu losowego oraz
        sieci MLP, CNN i LSTM,
  \item ocena każdego modelu (skuteczność oraz czas uczenia i działania),
  \item zebranie wyników w postaci tabel i wykresów.
\end{enumerate}

\subsection{Zbiór danych}

Głównym zbiorem danych jest \textbf{CSE-CIC-IDS2018} --- duży, publicznie
dostępny zbiór ruchu sieciowego zebranego na chmurowej infrastrukturze AWS przez
Canadian Institute for Cybersecurity~\cite{sharafaldin2018toward}. Zawiera on
zarówno ruch normalny, jak i różne typy ataków (m.in.\ łamanie haseł, DoS/DDoS,
ataki na aplikacje webowe, infiltrację oraz botnety). Każde połączenie opisane
jest kilkudziesięcioma cechami liczbowymi --- na przykład czasem trwania, liczbą
i rozmiarem przesłanych pakietów, tempem transmisji czy licznikami flag
sieciowych.

Pełny zbiór waży kilka gigabajtów, dlatego program wczytuje go w częściach i
losowo próbkuje, biorąc porcję danych z każdego dnia (każdy dzień zawiera inne
typy ataków). Dla zachowania pełnej powtarzalności --- gdyby ktoś chciał
uruchomić projekt bez pobierania wielkich plików --- przygotowano też
\emph{generator danych syntetycznych} o zbliżonej charakterystyce. Program
automatycznie korzysta z danych rzeczywistych, jeśli są dostępne, a w przeciwnym
razie z danych syntetycznych.

W opisywanym eksperymencie wykorzystano dane \texttt{\resDataSource}: łącznie
\resNTotal{} połączeń opisanych \resNFeatures{} cechami. Rozkład klas (widoczna
przewaga ruchu normalnego, typowa dla takich zadań) przedstawia
rys.~\ref{fig:class-dist}.

\begin{figure}[H]
\centering
\includegraphics[width=0.8\textwidth]{class_distribution.png}
\caption{Rozkład klas w danych. Ruch normalny wyraźnie dominuje nad atakami ---
to częste wyzwanie w wykrywaniu zagrożeń.}
\label{fig:class-dist}
\end{figure}

\subsection{Przetwarzanie wstępne danych}

Surowe dane wymagają porządków. Etap przygotowania obejmuje:
\begin{itemize}
  \item usunięcie kolumn nieprzydatnych do klasyfikacji (adresy IP, porty,
        znaczniki czasu) oraz kolumn o stałej wartości;
  \item poprawienie brakujących i błędnych wartości (np.\ wartości
        nieskończonych powstających przy dzieleniu);
  \item zamianę etykiet na prostą postać: ruch normalny lub atak (zachowano
        też informację o konkretnym typie ataku do celów poglądowych);
  \item ujednolicenie skali cech, tak aby duże liczby nie ,,zagłuszały''
        mniejszych --- co jest ważne zwłaszcza dla sieci neuronowych;
  \item podział danych na trzy części: uczącą, walidacyjną (do strojenia) i
        testową (do końcowej, uczciwej oceny).
\end{itemize}
Podział wykonano z ustalonym ziarnem losowości, dzięki czemu wyniki są
powtarzalne. Liczności zbiorów to: \resNTrain{} (uczący), \resNVal{}
(walidacyjny) i \resNTest{} (testowy) połączeń.

\subsection{Konfiguracja modeli}

Porównano pięć modeli o rosnącej złożoności:
\begin{itemize}
  \item \textbf{Regresja logistyczna} --- prosty model ,,liniowy'', traktowany
        jako punkt odniesienia: pokazuje, ile zyskujemy, sięgając po metody
        bardziej zaawansowane.
  \item \textbf{Las losowy} --- zespół 300 drzew decyzyjnych.
  \item \textbf{MLP} --- klasyczna sieć z dwiema warstwami ukrytymi.
  \item \textbf{CNN 1D} --- sieć konwolucyjna z dwiema warstwami filtrów.
  \item \textbf{LSTM} --- dwukierunkowa sieć rekurencyjna.
\end{itemize}
Aby uczciwie potraktować rzadkość ataków, w każdym modelu zastosowano
mechanizmy równoważące klasy (model jest ,,karany'' mocniej za przeoczenie
ataku). Sieci głębokie uczono z wykorzystaniem \emph{wczesnego zatrzymania} ---
nauka kończy się, gdy model przestaje się poprawiać na danych walidacyjnych, co
zapobiega przeuczeniu.

\subsection{Miary oceny i środowisko badawcze}

Skuteczność oceniono miarami opisanymi w rozdziale~\ref{sec:podstawy}
(dokładność, precyzja, czułość, F1, ROC-AUC i PR-AUC). Dodatkowo zmierzono
\emph{koszt}: czas potrzebny na wytrenowanie modelu, czas potrzebny na
sklasyfikowanie danych testowych oraz złożoność modelu (liczba parametrów lub
węzłów drzew). Całość zaimplementowano w języku \textbf{Python} przy użyciu
bibliotek \texttt{scikit-learn} (modele klasyczne) oraz \textbf{PyTorch} (sieci
głębokie); wykresy przygotowano w \texttt{matplotlib}. Obliczenia prowadzono na
zwykłym procesorze (CPU), co dobrze oddaje warunki ograniczonych zasobów i
uwypukla różnice w koszcie między metodami.
